% ==================================================
% فصل ترجمه مقاله Nature
% ==================================================

\section*{چکیده}
در این پروژه، مسئله‌ی پیش‌بینی حالت کامل یک سامانه‌ی چهار کیوبیتی بر اساس داده‌های خوانش کوادریچر مورد بررسی قرار گرفته است. داده‌ها شامل اندازه‌گیری‌های مختلط برای هر کیوبیت بوده و هدف، نگاشت هر نمونه‌ی اندازه‌گیری به یکی از ۱۶ حالت ممکن سامانه‌ی چهار کیوبیتی است. با وجود استفاده از روش‌های مختلف یادگیری ماشین کلاسیک، مهندسی ویژگی و کرنل‌های الهام‌گرفته از نگاشت‌های کوانتومی نظیر \lr{ZZ feature map}، دقت پیش‌بینی چندکلاسه در بهترین حالت به حدود ۳۸ تا ۳۹ درصد محدود شده است. در این گزارش، مسئله از دید فیزیکی و آماری تشریح شده، روش‌های پردازش داده و مدل‌سازی بررسی می‌شوند و در نهایت مسیرهای فیزیکی، آزمایشگاهی و الگوریتمی برای بهبود دقت مدل ارائه می‌گردد.

\section{مقدمه و شرح مسئله}
خوانش حالت کیوبیت‌ها یکی از چالش‌های اساسی در سامانه‌های محاسبات کوانتومی مبتنی بر پلتفرم‌هایی نظیر کیوبیت‌های ابررسانا یا یون‌های به‌دام‌افتاده است. در این سامانه‌ها، اطلاعات حالت کوانتومی به‌طور مستقیم قابل اندازه‌گیری نیست و از طریق فرآیندهای خوانش غیرمستقیم استخراج می‌شود. این فرآیندها معمولاً منجر به سیگنال‌هایی در فضای کوادریچر می‌شوند که به‌صورت نقاطی در صفحه‌ی مختلط \lr{(I,Q)} نمایش داده می‌شوند.

در این تحقیق، داده‌های مربوط به چهار کیوبیت مورد بررسی قرار گرفته‌اند. هر نمونه‌ی داده شامل چهار عدد مختلط است که هر یک متناظر با خوانش یک کیوبیت می‌باشد. برچسب هر نمونه یکی از ۱۶ حالت ممکن$ \ket{q_4 q_3 q_2 q_1}$ است. هدف اصلی، طراحی یک مدل طبقه‌بندی است که بتواند با دریافت این چهار خوانش مختلط، حالت کامل سامانه‌ی چهار کیوبیتی را پیش‌بینی کند.


\section{ساختار داده‌ها و تحلیل آماری اولیه}
هر مجموعه‌داده از ۱۶ فایل تشکیل شده است که هر فایل متناظر با یک حالت چهارکیوبیتی خاص می‌باشد. در هر فایل، برای هر کیوبیت ۱۰۰۰ نمونه‌ی اندازه‌گیری ثبت شده است. با ترکیب این داده‌ها، مجموعه‌داده‌ای شامل ۱۶۰۰۰ نمونه ایجاد شده است.

تحلیل‌های بصری اولیه شامل ترسیم نمودارهای پراکندگی در صفحه‌ی \lr{IQ}، هیستوگرام‌ها و توابع توزیع تجمعی انجام شد. این تحلیل‌ها نشان داد که توزیع داده‌های مربوط به حالات مختلف به‌شدت هم‌پوشان هستند. حتی با تقریب توزیع‌ها به شکل گاوسی و نمایش بیضی‌های کوواریانس، جدایش خطی یا حتی غیرخطی ساده بین کلاس‌ها امکان‌پذیر نبود. این رفتار با گزارش‌های تجربی رایج در خوانش کیوبیت‌ها هم‌خوانی دارد.

\section{مهندسی ویژگی و پیش‌پردازش داده}
در مرحله‌ی مهندسی ویژگی، دو نمایش اصلی برای داده‌های مختلط در نظر گرفته شد. در نمایش نخست، هر عدد مختلط به دو مولفه‌ی حقیقی و موهومی تجزیه شد که منجر به بردار ویژگی ۸بعدی می‌گردد. در نمایش دوم که به‌عنوان \lr{Angle Encoding} شناخته می‌شود، هر عدد مختلط به دامنه و فاز متناظر آن نگاشت شد. این نمایش از نظر فیزیکی معنادارتر بوده و به‌ویژه در ترکیب با کرنل‌های مبتنی بر توابع تناوبی مزیت دارد.

به‌منظور کاهش اثر داده‌های پرت و ناپایداری آماری، روش‌های مختلف نرمال‌سازی بررسی شدند. نتایج نشان داد که استفاده از \lr{RobustScaler} نسبت به \lr{StandardScaler} برای داده‌های نویزی آزمایشگاهی مناسب‌تر است.

\section{مدل‌سازی و الگوریتم‌های یادگیری ماشین}
مدل پایه‌ی مورد استفاده، ماشین بردار پشتیبان \lr{(SVM)} بود. در ابتدا، کرنل‌های استاندارد نظیر \lr{RBF} برای طبقه‌بندی تک‌کیوبیتی و چندکلاسه مورد استفاده قرار گرفتند. در حالت تک‌کیوبیتی، دقت‌های نسبتاً بالایی برای برخی کیوبیت‌ها به‌دست آمد، اما در حالت چندکلاسه، دقت به‌طور قابل توجهی کاهش یافت.

به‌منظور لحاظ کردن هم‌بستگی‌های بین ویژگی‌ها، یک کرنل سفارشی الهام‌گرفته از \lr{ZZ feature map} معرفی شد. این کرنل شامل ترم‌های تک‌ویژگی و ترم‌های جفتی بود که به‌صورت توابع کسینوسی تعریف می‌شدند. تنظیم پارامترهای این کرنل از طریق جستجوی شبکه‌ای انجام شد.

\section{نتایج و محدودیت‌ها}
بهترین دقت حاصل برای مسئله‌ی ۱۶کلاسه حدود ۳۸ تا ۳۹ درصد بود. این نتیجه نشان می‌دهد که اگرچه کرنل‌های الهام‌گرفته از کوانتوم می‌توانند اطلاعات بیشتری از ساختار داده‌ها استخراج کنند، اما هم‌پوشانی ذاتی توزیع‌ها و نویز اندازه‌گیری محدودیت اصلی باقی می‌ماند.

\section{مسیرهای پیشنهادی برای بهبود}
از منظر فیزیکی و آزمایشگاهی، بهبود تفکیک‌پذیری خوانش، کاهش \lr{cross-talk} و افزایش نسبت سیگنال به نویز می‌تواند نقش کلیدی ایفا کند. از منظر مدل‌سازی، استفاده از مدل‌های مولد، روش‌های \lr{ensemble} و مدل‌های گرافی پیشنهاد می‌شود. در نهایت، بهره‌گیری از الگوریتم‌های کوانتومی واقعی نظیر \lr{QSVM} یا طبقه‌بندهای واریاسیونی کوانتومی می‌تواند مسیر آینده‌ی این تحقیق باشد.


\section{تحلیل هم‌پوشانی ذاتی و وابستگی خوانش چندکیوبیتی}
به‌منظور تشخیص این‌که محدودیت دقت مدل ناشی از هم‌پوشانی ذاتی توزیع‌های خوانش است یا وابستگی و \lr{cross-talk} بین کیوبیت‌ها، یک تحلیل داده‌محور با الهام از معیارهای معرفی‌شده در مقاله‌ی
\lr{Multiplexed qubit readout quality metric beyond assignment fidelity}
انجام شد. تأکید این تحلیل بر استفاده‌ی صرف از داده‌های موجود خوانش بوده و هیچ‌گونه دسترسی به پارامترهای فیزیکی آزمایش فرض نشده است.

\subsection{جدایش‌پذیری ذاتی تک‌کیوبیتی}
در گام نخست، برای هر کیوبیت یک مدل طبقه‌بندی دودویی تنها بر اساس خوانش همان کیوبیت آموزش داده شد. دقت این مدل‌ها را می‌توان به‌عنوان سقف اطلاعاتی جدایش‌پذیری ذاتی هر کیوبیت در نظر گرفت. نتایج به‌صورت زیر به‌دست آمد:
\begin{center}
	\begin{tabular}{c c}
		\hline
		کیوبیت & دقت جدایش ذاتی \\
		\hline
		کیوبیت ۱ & 0.713 \\
		کیوبیت ۲ & 0.818 \\
		کیوبیت ۳ & 0.802 \\
		کیوبیت ۴ & 0.814 \\
		\hline
	\end{tabular}
\end{center}

این نتایج نشان می‌دهند که حتی در بهترین حالت و با استفاده از اطلاعات خوانش همان کیوبیت، خطای ذاتی قابل توجهی وجود دارد. اگر فرض شود خطاهای خوانش کیوبیت‌ها مستقل هستند، سقف دقت پیش‌بینی حالت چهارکیوبیتی را می‌توان به‌صورت تقریبی برابر حاصل‌ضرب دقت‌های تک‌کیوبیتی برآورد کرد که عددی در حدود ۳۸ درصد به‌دست می‌دهد. این مقدار با بهترین دقت مشاهده‌شده در مدل چندکلاسه هم‌خوانی دارد.

\subsection{تحلیل وابستگی بین‌کیوبیتی (Cross-talk)}
در گام دوم، وابستگی خوانش کیوبیت $i$ به حالت کیوبیت $j$ بررسی شد. برای این منظور، توزیع خوانش کیوبیت $i$ مشروط به $q_j=0$ و $q_j=1$ مقایسه گردید. معیارهایی نظیر جابه‌جایی میانگین و فاصله‌ی ماهالانوبیس بین این دو توزیع محاسبه شدند.

نتایج نشان دادند که عناصر روی قطر ماتریس وابستگی (یعنی $i=j$) به‌طور قابل توجهی بزرگ‌تر از عناصر خارج از قطر هستند، در حالی که وابستگی‌های بین‌کیوبیتی مقدار بسیار کوچکی دارند. اگرچه برخی از این وابستگی‌های ضعیف از نظر آماری معنی‌دار هستند، اما اندازه‌ی اثر آن‌ها از نظر فیزیکی ناچیز است. این موضوع نشان می‌دهد که \lr{cross-talk} بین خوانش‌ها نقش غالبی در محدودیت دقت مدل ایفا نمی‌کند.

\subsection{نتیجه‌گیری تحلیلی}
این تحلیل نشان می‌دهد که عامل اصلی محدودکننده‌ی دقت پیش‌بینی حالت چهارکیوبیتی، هم‌پوشانی ذاتی توزیع‌های خوانش تک‌کیوبیتی است، نه وابستگی شدید بین خوانش کیوبیت‌ها. بنابراین، افزایش قابل توجه دقت مدل تنها از طریق پیچیده‌تر کردن الگوریتم‌های یادگیری ماشین ممکن نیست و نیازمند بهبودهای فیزیکی در فرآیند خوانش و افزایش تفکیک‌پذیری اندازه‌گیری‌ها است.


\section{دیکودر گاوسی بیزی (تحلیل تمایز درجه دوم)}
به‌منظور به‌کارگیری یک رویکرد مدل‌محور متناسب با ماهیت آماری خوانش کیوبیت‌ها، روشی مبتنی بر مدل‌سازی مولد گاوسی و قانون بیز پیاده‌سازی و ارزیابی شد. این روش تحت عنوان \lr{Quadratic Discriminant Analysis (QDA)} یا دیکودر گاوسی بیزی شناخته می‌شود و در بسیاری از کاربردهای خوانش کیوبیت به‌عنوان یک روش استاندارد مورد استفاده قرار می‌گیرد.

\subsection{انگیزه و هدف}
در روش‌های تشخیصی معمول (مثل \lr{SVM})، به‌طور مستقیم توزیع شرطی برچسب به‌ازای ویژگی‌ها، یعنی $P(y \mid x)$، یادگیری می‌شود و مرز تصمیم به‌صورت هندسی برآورد می‌گردد. در مقابل، در رویکرد مولد گاوسی، توزیع داده به‌ازای هر حالت (کلاس) مدل می‌شود؛ یعنی $P(x \mid s)$ برای هر حالت $s$ از ۱۶ حالت ممکن. سپس با به‌کارگیری قانون بیز و با فرض prior یکنواخت (یا برآوردشده از داده)، حالت به‌عنوان
\begin{equation}
	\hat{s} = \operatorname*{arg\,max}_{s} \; P(x \mid s)\, P(s) = \operatorname*{arg\,max}_{s} \; P(x \mid s)
	\label{eq:bayes-decoder}
\end{equation}
انتخاب می‌شود (در صورت یکنواخت بودن $P(s)$). مزیت این رویکرد آن است که ساختار نویز و کوواریانس واقعی هر کلاس به‌طور صریح در مدل وارد می‌شود و مرز تصمیم به‌صورت بیضی‌وارهای (\lr{quadratic}) در فضای ویژگی شکل می‌گیرد، که با توزیع‌های گاوسی چندمتغیره هم‌خوانی بهتری دارد.

\subsection{فرض‌های مدل و فرم ریاضی}
فرض می‌کنیم برای هر حالت $s \in \{0,1,\ldots,15\}$، بردار ویژگی $x \in \mathbb{R}^8$ (شامل مولفه‌های حقیقی و موهومی چهار خوانش مختلط) از یک توزیع گاوسی چندمتغیره با میانگین $\mu_s$ و ماتریس کوواریانس $\Sigma_s$ تبعیت می‌کند:
\begin{equation}
	x \mid s \;\sim\; \mathcal{N}(\mu_s,\, \Sigma_s).
	\label{eq:gaussian-class}
\end{equation}
به‌جای کار با احتمال مستقیم، از لگاریتم احتمال (\lr{log-likelihood}) استفاده می‌کنیم که از نظر عددی پایدارتر است. برای یک نمونه $x$ و کلاس $s$:
\begin{equation}
	\log P(x \mid s) = -\frac{1}{2}(x - \mu_s)^\top \Sigma_s^{-1} (x - \mu_s) - \frac{1}{2} \log |\Sigma_s| + \mathrm{const}.
	\label{eq:log-likelihood}
\end{equation}
عبارت اول معادل منفی نصف فاصله‌ی ماهالانوبیس (مقیاس‌شده با $\Sigma_s$) و عبارت دوم مربوط به حجم بیضی کوواریانس است. در نهایت، کلاسی که مقدار $\log P(x \mid s) + \log P(s)$ را بیشینه می‌کند به‌عنوان پیش‌بینی انتخاب می‌شود؛ که با برآورد $P(s)$ از نسبت نمونه‌های هر کلاس در داده‌ی آموزش (\lr{prior}) انجام می‌شود.

\subsection{پارامترهای مدل و منظم‌سازی}
برای هر کلاس $s$، پارامترهای $\mu_s$ و $\Sigma_s$ به‌صورت تجربی از داده‌ی آموزش برآورد می‌شوند:
\begin{equation}
	\hat{\mu}_s = \frac{1}{n_s} \sum_{i:\, y_i = s} x_i, \qquad \hat{\Sigma}_s = \frac{1}{n_s-1} \sum_{i:\, y_i = s} (x_i - \hat{\mu}_s)(x_i - \hat{\mu}_s)^\top,
	\label{eq:mean-cov}
\end{equation}
که $n_s$ تعداد نمونه‌های کلاس $s$ است. برای جلوگیری از وارون‌ناپذیری عددی یا تکین شدن $\Sigma_s$ (به‌ویژه وقتی بعد ویژگی نسبت به تعداد نمونه زیاد است)، یک منظم‌ساز به صورت $\Sigma_s \leftarrow \Sigma_s + \varepsilon I$ با $\varepsilon = 10^{-6}$ اعمال شد.

در مرحله‌ی پیش‌بینی، به‌جای محاسبه‌ی صریح $\Sigma_s^{-1}$ (که از نظر عددی ناپایدار است)، از حل دستگاه خطی $\Sigma_s \, z = (x - \mu_s)$ با \lr{\texttt{np.linalg.solve}} استفاده شد؛ به‌طوری که ترم درجه‌دوم در \eqref{eq:log-likelihood} به صورت $-\frac{1}{2}(x - \mu_s)^\top z$ محاسبه می‌شود. همچنین برای $\log |\Sigma_s|$ از \lr{\texttt{np.linalg.slogdet}} استفاده شد تا از پایداری عددی در محدوده‌ی مقادیر بسیار کوچک یا بزرگ اطمینان حاصل شود.

\subsection{نرمال‌سازی ویژگی و Shrinkage}
در پیاده‌سازی اولیه، دقت مدل گاوسی بیزی حدود $38{,}67\%$ به‌دست آمد که با دقت مدل \lr{SVM} قابل مقایسه بود. برای بهبود، دو تغییر اعمال شد:
\begin{enumerate}
	\item \textbf{نرمال‌سازی ویژگی:} اعمال \lr{StandardScaler} روی داده‌های آموزش و آزمون تا هر ویژگی با میانگین صفر و واریانس واحد در فضای آموزش تبدیل شود. این کار فرض گاوسی را در فضای ویژگی متعادل‌تر می‌کند و پایداری عددی را افزایش می‌دهد.
	\item \textbf{Shrinkage کوواریانس:} در نسخه‌ی \lr{sklearn} از پارامتر \lr{\texttt{reg\_param}} استفاده شد که ماتریس کوواریانس هر کلاس را به سمت یک ماتریس مشترک «جمع» می‌کند (\lr{shrinkage}). این کار از overfitting و تکین شدن کوواریانس‌های کلاس‌هایی با نمونه کم جلوگیری می‌کند.
\end{enumerate}
جستجوی شبکه‌ای روی \lr{\texttt{reg\_param}} در بازه‌ی $[0,\, 0{,}5]$ انجام شد. بهترین دقت در حدود \lr{\texttt{reg\_param=0.30}} حاصل شد و دقت نهایی به حدود $39{,}20\%$ رسید؛ یعنی بهبودی در حد حدود نیم واحد درصد نسبت به حالت بدون shrinkage و نرمال‌سازی.

\subsection{تحلیل ماتریس آشفتگی}
ماتریس آشفتگی (\lr{confusion matrix}) مدل گاوسی بیزی نشان داد که بیشتر خطاها بین \lr{state}های «همسایه» رخ می‌دهند؛ یعنی حالت‌هایی که در فاصله‌ی همینگ ۱ از یکدیگر قرار دارند (فقط در یک کیوبیت با هم تفاوت دارند). این موضوع با فرض فیزیکی خطای خوانش تک‌کیوبیتی (\lr{bit flip}) سازگار است و نشان می‌دهد که مدل تا حدی ساختار خطا را درک می‌کند، اما به‌دلیل هم‌پوشانی زیاد توزیع‌ها در فضای هشت‌بعدی، جدایش کامل ۱۶ کلاس ممکن نیست.

\subsection{جمع‌بندی روش}
روش دیکودر گاوسی بیزی (\lr{QDA}) به‌صورت کامل پیاده‌سازی شد و با نرمال‌سازی و \lr{shrinkage} اندکی دقت را بهبود بخشید. با این حال، به‌دلیل هم‌پوشانی ذاتی توزیع‌های خوانش و محدودیت اطلاعات در هشت بعد برای ۱۶ کلاس، افزایش چشمگیر دقت تنها با این مدل قابل انتظار نبود. این نتیجه با تحلیل قبلی مبنی بر غلبه‌ی محدودیت فیزیکی خوانش تک‌کیوبیتی بر محدودیت الگوریتم هم‌خوانی دارد. برای گزارش نهایی، کدهای مربوط به بارگذاری داده، تبدیل به ویژگی حقیقی هشت‌بعدی، برآورد میانگین و کوواریانس هر کلاس، تابع \lr{log-likelihood} با \lr{\texttt{solve}} و \lr{\texttt{slogdet}}، و پیش‌بینی با بیشینه‌سازی امتیاز بیزی در نوت‌بوک پروژه موجود است.


% ==================================================
\section{مدل مخلوط گاوسی (Gaussian Mixture Model) برای دیکودینگ چندحالته}
% ==================================================

\subsection{انگیزه و چارچوب روش}

در ادامه‌ی رویکرد مولد گاوسی تک‌مدلی (QDA)، این فرض مورد بازبینی قرار گرفت که توزیع داده‌های خوانش هر حالت الزاماً تک‌قله‌ای (Unimodal) نیست. در سامانه‌های خوانش چندکیوبیتی، عواملی نظیر نویز تقویت‌کننده، تغییرات فاز، دریفت زمانی و نوسانات محیطی می‌توانند باعث ایجاد ساختارهای چندمدی در توزیع خوانش شوند. در چنین شرایطی، تقریب هر کلاس با یک توزیع گاوسی منفرد ممکن است ساده‌سازی بیش از حد باشد.

به‌منظور مدل‌سازی دقیق‌تر ساختار آماری داده‌ها، برای هر کلاس $s$ از یک مدل مخلوط گاوسی استفاده شد:

\begin{equation}
	P(x \mid s) = \sum_{k=1}^{K_s} \pi_{s,k}\,
	\mathcal{N}(x \mid \mu_{s,k}, \Sigma_{s,k}),
	\label{eq:gmm-class}
\end{equation}

که در آن $K_s$ تعداد مؤلفه‌های گاوسی کلاس $s$، ضرایب $\pi_{s,k}$ وزن‌های مخلوط و $\mu_{s,k}$ و $\Sigma_{s,k}$ به‌ترتیب میانگین و کوواریانس مؤلفه‌ی $k$ام هستند. پارامترها با استفاده از الگوریتم بیشینه‌سازی امید ریاضی (\lr{Expectation-Maximization}) برآورد شدند.

پس از برآورد مدل برای هر کلاس، پیش‌بینی مشابه حالت QDA از طریق بیشینه‌سازی لگاریتم احتمال بیزی انجام شد:

\begin{equation}
	\hat{s} = \operatorname*{arg\,max}_{s}
	\left[ \log P(x \mid s) + \log P(s) \right].
\end{equation}

\subsection{تنظیم تعداد مؤلفه‌ها}

برای بررسی تأثیر پیچیدگی مدل، تعداد مؤلفه‌های مخلوط برای همه‌ی کلاس‌ها به‌صورت یکسان در بازه‌ی
\[
K \in \{1,2,3,4,5\}
\]
آزمایش شد. همچنین یک راهبرد تطبیقی بر اساس معیار اطلاعات بیزی (\lr{BIC}) به‌کار گرفته شد که در آن برای هر کلاس، مقدار بهینه‌ی $K_s$ در بازه‌ی فوق انتخاب می‌شد.

\subsection{نتایج عددی}

دقت مدل برای مقادیر مختلف تعداد مؤلفه‌ها به‌صورت زیر به‌دست آمد:

\begin{center}
	\begin{tabular}{c c}
		\hline
		تعداد مؤلفه‌ها & دقت \\
		\hline
		$K=1$ & $38.70\%$ \\
		$K=2$ & $37.60\%$ \\
		$K=3$ & $37.45\%$ \\
		$K=4$ & $36.08\%$ \\
		$K=5$ & $35.65\%$ \\
		\hline
	\end{tabular}
\end{center}

در حالت انتخاب تطبیقی مبتنی بر BIC نیز دقت نهایی برابر با $38.25\%$ به‌دست آمد. در این حالت، برخی کلاس‌ها به‌صورت بهینه با دو مؤلفه و برخی با یک مؤلفه مدل شدند، اما این افزایش انعطاف‌پذیری منجر به بهبود معنادار دقت نشد.

\subsection{تحلیل نتایج}

نتایج نشان می‌دهد که افزایش تعداد مؤلفه‌های گاوسی نه‌تنها باعث بهبود دقت نشد، بلکه در مقادیر بزرگ‌تر $K$ موجب کاهش عملکرد نیز گردید. این رفتار را می‌توان از چند منظر تفسیر کرد:

\begin{enumerate}
	\item \textbf{ماهیت تقریباً تک‌مدی توزیع‌ها:} تحلیل‌های بصری پیشین نشان می‌دهد که توزیع خوانش هر حالت در فضای ویژگی هشت‌بعدی تا حد زیادی تک‌مدی است. بنابراین، مدل تک‌گاوسی (QDA) تقریب مناسبی ارائه می‌دهد و افزایش پیچیدگی مدل ضروری نیست.
	
	\item \textbf{بیش‌برازش آماری:} با افزایش $K$، تعداد پارامترهای مدل به‌طور چشمگیری افزایش می‌یابد (به‌ویژه در حضور کوواریانس کامل). این امر می‌تواند به بیش‌برازش داده‌ی آموزش و کاهش تعمیم‌پذیری در داده‌ی آزمون منجر شود.
	
	\item \textbf{غلبه‌ی هم‌پوشانی بین‌کلاسی:} مهم‌ترین عامل محدودکننده‌ی دقت، هم‌پوشانی ذاتی بین توزیع‌های کلاس‌های مختلف است. افزودن مؤلفه‌های بیشتر درون هر کلاس نمی‌تواند این هم‌پوشانی بین‌کلاسی را برطرف کند، زیرا منشأ آن ناشی از محدودیت تفکیک‌پذیری خوانش تک‌کیوبیتی است.
\end{enumerate}

\subsection{جمع‌بندی}

مدل مخلوط گاوسی به‌عنوان تعمیمی بر دیکودر گاوسی بیزی پیاده‌سازی و ارزیابی شد. نتایج نشان داد که بهترین عملکرد زمانی حاصل می‌شود که هر کلاس با یک مؤلفه‌ی گاوسی مدل شود (معادل QDA). افزایش تعداد مؤلفه‌ها نه‌تنها بهبود قابل توجهی ایجاد نکرد، بلکه در اغلب موارد موجب افت دقت گردید.

این نتیجه تأیید می‌کند که محدودیت عملکرد مدل بیش از آن‌که ناشی از ساده‌سازی توزیع درون‌کلاسی باشد، ناشی از هم‌پوشانی ذاتی توزیع‌های خوانش در بین کلاس‌های مختلف است. بنابراین، افزایش پیچیدگی مدل مولد در سطح کلاس نمی‌تواند به‌تنهایی منجر به جهش قابل توجه در دقت پیش‌بینی شود.

\section{جمع‌بندی}
در این گزارش، مسئله‌ی پیش‌بینی حالت کامل چهار کیوبیت با استفاده از داده‌های خوانش کوادریچر بررسی شد. نتایج نشان داد که محدودیت دقت مدل بیش از آنکه ناشی از ضعف الگوریتم‌های یادگیری ماشین باشد، ریشه در ماهیت فیزیکی فرآیند اندازه‌گیری دارد. مسیرهای پیشنهادی ارائه‌شده می‌توانند مبنایی برای ادامه‌ی تحقیق در حوزه‌های آزمایشگاهی و الگوریتمی باشند.